% Agradecimentos (texto do aluno: 05_propostas_de_texto/agradecimentos.md no hub).
\begin{otherlanguage}{brazilian}
\chapter*{Agradecimentos}

Meus agradecimentos são curtos. Antes de tudo, vale dizer pelo que agradeço. Agradeço pelas oportunidades incríveis e pelas pessoas incríveis que tive ao meu redor. Agradeço por ter tido a chance de me testar e de mostrar o meu melhor. Agradeço a quem me apoiou e a mim mesmo, por ter tido coragem. Agradeço por todos os momentos difíceis e desagradáveis, por todas as dores e alegrias.

O que esses muitos anos me mostraram, e o que esta conclusão de curso sintetiza, é que tenho motivos para agradecer, porque a vida me ensinou o que é sentir, nas vitórias e nas derrotas. Dito isso, agradeço.

Como manda a tradição acadêmica, começo pelos professores que tive na vida, tanto os de quem gostei muito quanto os de quem nem tanto. A profissão de professor é muito sofrida e merece reconhecimento. Obrigado a todos que me ensinaram. Eu nunca teria chegado aonde estou sem vocês.

Obrigado às escolas que me apoiaram, que me deram bolsa e que acreditaram em mim. Obrigado à minha família: meus tios, minhas tias, minha avó e meu avô. Muito obrigado aos meus irmãos. Obrigado, Naty, minha irmãzinha linda, que está sempre torcendo por mim e que me diverte muito. Sou muito feliz por ter você como a pessoa geneticamente mais próxima de mim, porque te acho uma pessoa incrível.

Muito obrigado, pai, por me ensinar a ser resiliente, a sofrer e a lidar com o sofrimento, assim como por me apoiar nos meus sonhos mais malucos. Obrigado aos meus cachorros Bob e Kelsen, que me fizeram companhia na minha infância e adolescência.

E muitíssimo obrigado a todos os meus amigos por esta jornada. Nunca fui uma pessoa disciplinada, nunca tive muita rotina e nunca fui muito certinho. Eu teria sido mais um potencial desperdiçado pela procrastinação se não fossem as pessoas ao meu lado. Se hoje sou alguém, é porque as pessoas à minha volta valem dez vezes mais do que eu. Muito obrigado aos amigos do cursinho e do colégio, que hoje se tornaram amigos de ITA. Espero, do fundo do coração, que a nossa amizade seja para a vida inteira. Nestes tempos em que poucos valorizam os laços e tudo é tão frágil, quero ser um homem feliz cultivando essa amizade com vocês, para sempre.

Vocês são a minha família aqui, tão longe da minha família de sangue. Contem comigo sempre que precisarem. Espero ser uma boa família para vocês também. Não é fácil ficar longe. Não é fácil passar por tudo isso: pelas incertezas, pelas dificuldades, pelos momentos bons, pelo sofrimento e pela felicidade, às vezes muito sozinho. Que eu nunca deixe vocês se sentirem sozinhos. Espero cumprir esse papel como amigo.

Muito obrigado às empresas que me acolheram: Goldman Sachs e Enter, pelas oportunidades. Sendo honesto, sou muito sortudo por ter tido trabalhos em que me diverti de verdade e de que gostei muito.

Muito obrigado, em especial, aos meus dois primeiros chefes e companheiros de equipe, Gabriel Liberman e André Fioravanti. Espero um dia ser um décimo do profissional que vocês são. Obrigado pela paciência, por me ensinarem, por me fazerem melhor e por me tornarem o profissional que sou hoje. Devo tudo a vocês. Vocês não são apenas exemplos de profissionais. São exemplos de pessoas, de adultos, de ética, de responsabilidade e de competência. São os exemplos que quero seguir.

Muito obrigado também ao ITA, que, apesar dos pesares, me deu casa, comida e, espero, um diploma.

Obrigado a todas as garotas que já amei. Isso faz parte da vida, e que venha muito mais.

Me desculpem todos por às vezes ser meio exagerado ou intenso, mas acho que sou assim. Então, obrigado a mim mesmo por ser exagerado: por ter exagerado nos meus sonhos e desejado não só passar, mas passar em primeiro lugar; por ter desejado não só trabalhar, mas trabalhar nos melhores lugares. Obrigado a mim mesmo por ter exagerado no amor, nas paixões, nas amizades, na carência, na curiosidade, nos estudos, e em tudo aquilo que nos faz humanos. E também por nunca ter desistido nos momentos em que não alcancei esse exagero. Eu definitivamente não sou o melhor, mas eu tento, e tentar é divertido.

Por último, mas não menos importante, obrigado à pessoa mais importante da minha vida, e, a meu ver, uma das mais importantes do mundo. Obrigado, mãe. Você é para mim um exemplo de mulher, de homem, de pessoa, de inteligência, de caráter, de humildade e de bondade.

Às vezes acredito, de verdade, que você é uma santa. Não no sentido cego da palavra, mas no sentido de que talvez more em você o ápice da bondade e da benevolência. E não só isso: você é uma das pessoas mais inteligentes que já conheci, e isso não é pouca coisa.

Obrigado por me fazer 50\% você. Quem me dera fossem 100\%. Você me deu uma infância perfeita, uma adolescência perfeita e, agora, uma juventude adulta perfeita. Você não é a melhor mãe que eu poderia ter escolhido. Você simplesmente é a melhor. Te amo muito.

Acho que é isso. Esses foram os meus agradecimentos, e espero não ter esquecido ninguém. Desculpem o texto, talvez um pouco melodramático, mas acho que um pouco de melodrama às vezes faz parte: um pouco de exagero, um pouco de sentimentalismo, um pouco de esquenta e esfria, de aperta e afrouxa, de sossego e de desinquietação. Como diria Guimarães Rosa, o que importa é ter coragem.

\vspace{1cm}
\begin{flushright}\begin{minipage}{0.85\textwidth}\raggedleft
\foreignlanguage{english}{\itshape ``I'm sad, but at the same time, I'm really happy that something could make me feel that sad. It's like\ldots{} it makes me feel alive, you know? It makes me feel human.''}\par
(Butters Stotch, \emph{South Park})
\end{minipage}\end{flushright}
\end{otherlanguage}
\clearpage
